% ==============================================================================
% Fast Main-Text Chinese Review Template (XeLaTeX + ctexart)
% Mirrors English INFORMS main text 1-to-1 (excludes Online Appendix).
% Compile with: xelatex -interaction=nonstopmode -halt-on-error main_zh.tex
% ==============================================================================

% Synced from ../main.tex @ BASELINE_COMMIT
\documentclass[UTF8, 11pt, a4paper]{ctexart}

% --- Generous margins for iPad/PDF annotation ---
\usepackage[top=2.4cm, bottom=2.4cm, left=2.5cm, right=2.5cm]{geometry}
\usepackage{setspace}
\onehalfspacing

% --- Math, Theorems, Tables, Graphics, Algorithms ---
\usepackage{amsmath,amssymb,amsfonts,mathtools,bm}
\usepackage{amsthm}
\usepackage{booktabs,threeparttable,multirow}
\usepackage{graphicx}
\graphicspath{{../}{../figures/}{../figs/}}
\usepackage{algorithm}
\usepackage{algpseudocode}
\usepackage{xcolor}
\usepackage[colorlinks=true,linkcolor=blue!60!black,citecolor=blue!60!black,urlcolor=blue!60!black]{hyperref}

% --- Theorem Environments (Keep numbering parallel to English main.tex) ---
\theoremstyle{plain}
\newtheorem{theorem}{定理}
\newtheorem{proposition}{命题}
\newtheorem{lemma}{引理}
\newtheorem{corollary}{推论}

\theoremstyle{definition}
\newtheorem{definition}{定义}
\newtheorem{assumption}{假设}

\theoremstyle{remark}
\newtheorem{remark}{注}

% --- Convenience Macros matching English template ---
\newcommand{\E}{\mathbb{E}}
\newcommand{\R}{\mathbb{R}}

\title{\textbf{论文中文快速审阅版标题（与英文主文一一对应）}}
\author{审阅模式（已略去 Online Appendix 证明与补充实验）}
\date{}

\begin{document}
\maketitle

% [EN-SYNC: main.tex:abstract:p1]
\begin{abstract}
本文研究动态定价与库存协同优化问题。在需求存在周期性波动与不确定性的情形下，传统静态策略往往难以兼顾收益与缺货风险。我们构建了一个多周期随机动态规划模型，证明了最优策略具有阈值结构，并提出了一种基于拉格朗日松弛的可计算近似算法。数值结果表明，相比传统启发式基准策略，所提出的自适应阈值策略在保持计算可控的同时，显著缩小了收益差距。
\end{abstract}

\noindent\textbf{关键词：} 动态定价，库存控制，阈值策略，随机优化

% [EN-SYNC: main.tex:sec1:p1]
\section{引言 (Introduction)}
\label{sec:intro}

在零售与平台运营中，决策者需要在需求尚未完全明朗前同时设定库存补货量与销售价格。当供应提前期较长且市场需求呈现显著波动时，静态定价策略无法及时响应库存状态的变化，进而导致缺货损失或期末积压。

% [EN-SYNC: main.tex:sec2:p1]
\section{模型设定 (Model)}
\label{sec:model}

考虑一个离散有限周期 $t \in \{1, \dots, T\}$ 的运营系统。在每一期期初，决策者观察当前库存状态 $x_t$ 后选择当期补货量与零售价格，随后随机需求实现并产生相应收益或持有成本。

% [EN-SYNC: main.tex:prop1]
\begin{proposition}[阈值最优性]
\label{prop:threshold}
当期需求分布满足单调失效率条件时，最优定价与补货策略在每一期均保持单阈值结构。
\end{proposition}

\vspace{1em}
\noindent\textit{注：Online Appendix（详细证明、辅助引理及补充数值实验）已按快速审阅模式略去。}

\end{document}
